\chapter{Discusión}
\label{cap:discusion}

\section{Evaluación de las hipótesis}

\subsection*{H1: disparidades entre entidades federativas}

\textbf{Sostenida.} Un modelo con AUROC de \aurocGlobal{} presenta una brecha de
sensibilidad de \gapTPRentidad{} y una brecha de AUROC de \gapAUROCentidad{} entre
entidades. La disparidad persiste al restringir el análisis a adultos de 60 años o
más, donde el rango de AUROC va de \sesentaPeorAUROC{} a \sesentaMejorAUROC, lo
que descarta que sea un reflejo de la composición etaria de los estados.

Resiste además dos pruebas de robustez que podrían haberla desmentido. En el
subconjunto hospitalizado --la población en la que el desenlace es posible-- no
solo persiste sino que la brecha de AUROC casi se triplica, hasta \hospGapAUROC{}
(sección~\ref{sec:sensibilidad_res}). Y no se debe a la composición del
entrenamiento: un modelo que nunca vio una entidad rinde en ella casi igual que
uno que sí la vio (sección~\ref{sec:loeo_res}).

La formulación original hablaba de disparidades «estadísticamente
significativas», y esa cualificación ahora sí puede sostenerse. Ninguna de las
\boP{} reasignaciones al azar de la etiqueta de entidad produjo una brecha tan
grande como la observada, ni en AUROC ni en sensibilidad ($p=\pAuroc$ y
$p=\pTpr$, el mínimo alcanzable con ese número de permutaciones), y
\paresPct\,\% de los pares de entidades tienen intervalos de AUROC disjuntos
(sección~\ref{sec:incertidumbre_res}).

El contraste obliga además a matizar la magnitud. La brecha nula no es cero: con
32 grupos desiguales, el azar produce una separación mediana de \nulaTpr{} en
sensibilidad. El exceso atribuible a la entidad es de unos \excesoTpr, no los
\gapTPRentidad{} completos. Sigue siendo una disparidad grande, pero conviene
reportarla contra la referencia correcta y no contra cero.

\subsection*{H2: asociación con la infraestructura o marginación}

\textbf{No sostenida.} El índice de marginación de CONAPO 2020 no muestra
asociación detectable con el desempeño por entidad ($r=\corrMargAurocr$,
$p=\corrMargAurocp$), y el resultado se sostiene al agregar por grado de
marginación, donde una prueba de Kruskal--Wallis tampoco detecta diferencia
($H=\kruskalH$, $p=\kruskalP$). El grupo con menor AUROC medio es, de hecho, el
de marginación \emph{muy baja}: si hay una tendencia, apunta en dirección
opuesta a la hipótesis.

Cabe una advertencia sobre el alcance de esta conclusión. La ausencia de
asociación con un índice compuesto no equivale a ausencia de mecanismo
estructural: la marginación socioeconómica agregada por entidad es un
\emph{proxy} grueso de lo que importaría aquí, que es la capacidad de registro y
atención de las unidades médicas. \ref{h2} queda refutada en la forma en que se
operacionalizó, no en su intuición de fondo.

El único correlato que alcanza significación convencional es la tasa de mortalidad
de la entidad ($r=\corrAurocTasar$, $p=\corrAurocTasap$): el modelo discrimina
peor donde más gente muere. Conviene no sobrevalorarlo: se examinaron tres
características por dos métricas y dos estadísticos, de modo que un único
$p$ apenas por debajo de 0.05 en una docena de contrastes no sobrevive a ninguna
corrección por multiplicidad. Se reporta como señal a investigar, no como
hallazgo confirmatorio. Esta asociación admite al menos dos lecturas que el diseño actual no puede
separar. Puede reflejar un mecanismo de \emph{atención}: en entidades con sistemas
más presionados, la muerte depende más de factores no capturados por el modelo
--disponibilidad de cama, oportunidad del traslado, saturación-- y menos del perfil
clínico inicial, de modo que hay un techo real a lo que la edad y las
comorbilidades pueden predecir. Y puede reflejar un mecanismo de \emph{registro}:
una mortalidad observada más alta puede indicar subdetección de casos leves, lo que
cambia la composición de la cohorte y con ella la dificultad de la tarea.
Distinguir ambas requiere indicadores de infraestructura por unidad médica, no por
entidad.

\subsection*{H3: mitigación con pérdida acotada de desempeño}

\textbf{Sostenida, con una precisión importante.} El post-procesamiento por
umbrales de entidad reduce la brecha de sensibilidad de \gapBase{} a \gapPost{}
--\reduccionPost\,\%-- sin pérdida de desempeño agregado: la exactitud balanceada
pasa de \accBase{} a \accPost. La frontera de Pareto muestra que el desempeño se
mantiene en la banda \paretoAccMin--\paretoAccMax{} en todo su recorrido, de modo que la pérdida es
inferior a cualquier umbral razonable que se hubiera fijado de antemano.

Hay, sin embargo, un costo que la exactitud balanceada no captura y que conviene
enunciar: al igualar la sensibilidad entre entidades, la brecha de \emph{tasa de
falsos positivos} empeora, de \postGapFPRbase{} a \postGapFPR. No es un defecto
del método sino exactamente lo que predice el teorema de imposibilidad
(capítulo~\ref{cap:marco}): con prevalencias distintas entre grupos, igualar una
de las dos tasas de \emph{equalized odds} desiguala la otra. Lo que el
post-procesamiento hace no es eliminar la inequidad, sino \emph{elegir} cuál de
las dos se reparte por igual --y elegir la sensibilidad, en un instrumento cuyo
costo de omitir una defunción excede al de una alarma innecesaria, es defendible.
Pero es una elección normativa que debe declararse, no un óptimo técnico.

La segunda precisión es que la mitigación \emph{no elimina} la disparidad. Una
brecha residual de \gapPost{} sigue siendo clínicamente relevante. Y el hallazgo
metodológico que acompaña al resultado es que la cifra de eliminación casi total
que se obtiene ajustando umbrales sobre el propio conjunto de prueba
(\gapPostOpt) es un artefacto de ajuste en muestra. Un trabajo que no separase el
bloque de calibración concluiría que el problema queda resuelto, cuando la brecha
desplegable es \factorOptimismo{} veces mayor.

\subsection*{H4: transportabilidad a una fuente independiente}

\textbf{Sin poner a prueba, con evidencia parcial en la dirección contraria a lo
esperado.} La validación externa en la base de Egresos Hospitalarios de la DGIS no
se ha ejecutado, de modo que \ref{h4} tal como está redactada --transferencia a
una \emph{fuente independiente}-- no queda ni sostenida ni refutada.

Sí se puso a prueba la generalización geográfica dentro de la misma fuente
(sección~\ref{sec:loeo_res}), y el resultado va en contra de lo que \ref{h4}
anticipaba: al excluir por completo una entidad del entrenamiento, el desempeño
sobre ella apenas cambia (pérdida mediana \loeoMediana) y las disparidades no se
acentúan, sino que se mantienen o se reducen. Esto acota la hipótesis: si las
disparidades llegaran a agravarse en la base DGIS, no sería por no haber visto la
entidad, sino por el cambio de fuente, de definición de caso o de periodo.

La ejecución futura de \ref{h4} tiene además un prerrequisito que este trabajo
identificó y conviene dejar escrito. SAEH es una base de \emph{egresos}: su
población son pacientes hospitalizados, con una prevalencia del desenlace cercana
a \hospPrev{} frente a \cohorteTasaDef\,\% en esta cohorte. Aplicar sobre ella el
modelo principal produciría una caída de desempeño atribuible al cambio de
espectro y no a la transportabilidad. La validación externa debe hacerse con el
modelo restringido a hospitalizados de la sección~\ref{sec:sensibilidad_res}, no
con el de la cohorte completa.

\section{Asimetría entre familias de mitigación}

El resultado comparativo tiene, a nuestro juicio, valor propio. Las tres familias
de mitigación se presentan habitualmente como alternativas de eficacia
comparable, a elegir según restricciones operativas. En este problema la
diferencia es cualitativa: el pre-procesamiento no mueve la brecha, el
in-procesamiento arroja un resultado tan disperso que no permite concluir nada
--su rango entre realizaciones, $[\inMin, \inMax]$, contiene a la línea de
base-- y solo el post-procesamiento la reduce, a la mitad y sin costo agregado.

La explicación está en la naturaleza del sesgo. La disparidad no proviene de un
ordenamiento defectuoso: el AUROC por entidad es alto en todas
(\entPeorAUROC--\entMejorAUROC), y el modelo está bien calibrado en cada una
(error máximo \errCalEntidad). Proviene de aplicar un \emph{corte único} a
distribuciones de riesgo que difieren entre entidades. Cuando el defecto está en
la regla de decisión y no en la función de puntuación, la intervención que corrige
la regla de decisión es la que funciona. Reponderar los datos de entrenamiento
ataca un problema de representación que aquí no es el que causa la brecha.

Esto sugiere una heurística transferible: antes de elegir la familia de
mitigación, conviene diagnosticar si la inequidad reside en el ordenamiento
--brechas de AUROC grandes-- o en el umbral --AUROC similar con brechas de TPR
grandes. Aquí la brecha de sensibilidad (\gapTPRentidad) es \razonBrechas{} veces
la de AUROC (\gapAUROCentidad), lo que apuntaba al umbral desde el diagnóstico.
Ese diagnóstico, sin embargo, es específico de la cohorte completa. Entre los
pacientes hospitalizados la relación se invierte --la brecha de AUROC sube a
\hospGapAUROC{} y la de sensibilidad baja a \hospGapTPR-- y en el subgrupo de 60
años o más la dispersión del AUROC llega a \sesentaRango. La heurística, por
tanto, no es «diagnostíquese una vez», sino «diagnostíquese en la población sobre
la que se va a desplegar»: en la población en riesgo, una mitigación basada solo
en umbrales sería insuficiente.

\section{Implicaciones para el despliegue}

\begin{enumerate}[leftmargin=*]
  \item \textbf{Un umbral nacional único no es una decisión neutral.} Es una
    decisión que redistribuye sensibilidad entre entidades. Conviene precisar en qué
    dirección: la sensibilidad por entidad \emph{no} está asociada a su mortalidad
    ($r=\corrTprTasar$, $p=\corrTprTasap$), de modo que el reparto no penaliza
    sistemáticamente a los estados con más defunciones; lo que sí decae en ellos es
    la discriminación ($r=\corrAurocTasar$, $p=\corrAurocTasap$). La inequidad de
    sensibilidad no sigue un eje interpretable, y eso la hace más difícil de
    justificar, no menos: la neutralidad aparente de «el mismo umbral para todos»
    reparte la protección de forma esencialmente arbitraria entre estados.
  \item \textbf{La frontera de Pareto es el entregable adecuado para el
    regulador.} La elección de $\alpha$ es normativa. El análisis técnico puede
    mostrar que no cuesta desempeño agregado; no puede decidir cuánta
    diferenciación por entidad es aceptable en términos de equidad horizontal
    entre ciudadanos.
  \item \textbf{Umbrales por entidad plantean una objeción legítima.} Usar la
    residencia en el momento de la decisión implica tratar de forma distinta a dos
    pacientes con idéntico perfil clínico. La justificación es que el trato
    formalmente idéntico ya produce resultados desiguales; pero es un argumento que
    debe hacerse explícito ante un comité de ética, no asumirse.
  \item \textbf{La calibración debe verificarse antes de auditar equidad.} Como
    muestra el capítulo~\ref{cap:resultados}, auditar sobre puntajes descalibrados
    habría producido una conclusión invertida sobre la calibración por grupo.
  \item \textbf{Humano en el circuito.} Con una tasa de selección de \selGlobal{} y un
    valor predictivo positivo modesto en el punto de operación, el instrumento es
    de priorización y no de decisión automatizada.
\end{enumerate}

\section{Relación con la literatura}

El modelo de referencia coincide con la literatura sobre esta
base~\parencite{becerra2022,mendez2024} en las variables que resultan más
informativas --neumonía y edad-- lo que respalda que el sujeto de la auditoría sea
representativo y no un modelo debilitado. Una comparación numérica directa del
desempeño no es posible, porque esos trabajos reportan exactitud y $F_1$ sobre
umbrales no declarados y esta tesis reporta AUROC y sensibilidad en un punto de
operación explícito.

Respecto del trabajo que estudia los factores geográficos como
predictores~\parencite{geograficos2024}, los resultados de esta tesis son
complementarios y no contradictorios. Que la entidad sea informativa como
predictor y que el modelo sea desigualmente preciso entre entidades son hechos
distintos; el segundo es el que importa para el despliegue responsable, y es el
que no se había medido. Es de notar que el modelo auditado aquí \emph{no} usa la
entidad como predictor, y aun así su desempeño depende de ella: la inequidad
geográfica no requiere que la geografía entre al modelo.

Frente a la literatura de equidad en salud anclada en otros
contextos~\parencite{nejadshamsi2026,riccilara2022}, la aportación es el traslado del
marco a un atributo sensible --la entidad federativa-- con significado
institucional propio en México, donde la desigualdad entre estados es un rasgo
estructural del sistema y no una variable de composición poblacional. Que el
índice de marginación no explique la brecha no debilita esa elección: confirma
que la entidad captura algo que los agregados socioeconómicos estatales no
miden, y esa es precisamente la razón para auditarla como atributo propio.

\section{Limitaciones}
\label{sec:limitaciones}

\begin{enumerate}[leftmargin=*]
  \item \textbf{Sin validación externa.} Todo el análisis vive en una fuente y un
    año. La generalización geográfica interna sí se evaluó
    (sección~\ref{sec:loeo_res}) y resultó buena, pero no sustituye a una segunda
    fuente.
  \item \textbf{El desenlace vive en el \fracHosp\,\% de la cohorte.} Fallecer
    implica estar registrado como hospitalizado, de modo que parte de la
    discriminación del modelo principal es separar ambulatorios de hospitalizados.
    Las cifras comparables con modelos clínicos de mortalidad son las de la
    sección~\ref{sec:sensibilidad_res}, no las de la cohorte completa.
  \item \textbf{Incertidumbre acotada solo en el eje principal.} La auditoría
    por entidad se acompaña de intervalos de confianza y prueba de permutación,
    pero los subgrupos interseccionales y los ejes de sexo y edad se reportan como
    estimaciones puntuales. El ordenamiento fino de la tabla~\ref{tab:entidad}
    tampoco debe leerse posición por posición: solo \paresPct\,\% de los pares
    tienen intervalos disjuntos.
  \item \textbf{Marginación medida a escala estatal.} El índice de CONAPO está
    verificado, pero agrega condiciones socioeconómicas por entidad. El mecanismo
    que \ref{h2} postula --capacidad de atención-- opera a escala de unidad médica,
    de modo que la ausencia de asociación acota la hipótesis pero no la cierra.
  \item \textbf{Un solo año y un solo desenlace.} No se modela hospitalización ni
    se evalúa estabilidad temporal entre olas epidémicas, cuya composición de casos
    y presión hospitalaria difieren de forma conocida.
  \item \textbf{Sesgo de verificación en la cohorte.} Restringir a casos
    confirmados condiciona a haber sido detectado y probado, y la propensión a
    detectar varía entre entidades. Parte de la disparidad observada puede
    originarse en el proceso de captura y no en el modelo. Este es, a nuestro
    juicio, el supuesto más frágil del diseño y no puede corregirse dentro de esta
    fuente.
  \item \textbf{Dependencia de una variable.} Con \texttt{NEUMONIA} concentrando
    la importancia, un cambio en la práctica de registro entre entidades se
    traduciría directamente en desempeño desigual. Falta el análisis de
    sensibilidad que lo descarte.
  \item \textbf{Validación por partición aleatoria.} No se realizó validación
    temporal ni dejando-una-entidad-fuera, que es el diseño natural para evaluar
    generalización geográfica.
\end{enumerate}
